% =====================================================================
\chapter{Soluciones de los ejercicios}
\label{ap:soluciones}
% =====================================================================

Cada solución se presenta en cinco partes: datos, planteamiento, desarrollo, resultado y comprobación. En los pasos intermedios se redondea a cuatro cifras significativas salvo que se indique otra cosa; los resultados exactos se dan como fracción cuando existe. Las memorias se expresan con $1\,\mathrm{GB}=10^{9}$ bytes.

% ---------------------------------------------------------------------
\section{Tema 1}

\begin{solucion}{eje:1-tasa}
\planteamiento el cuerpo parte de $\theta_0$, que ya es una buena solución; la tasa de aprendizaje controla cuánto se aleja de ella (\cref{obs:implicita}).
\pasos
\begin{enumerate}[label=\arabic*.]
  \item Con tasas pequeñas, el cuerpo se adapta poco a poco y la cabeza, que parte de cero, aprende a usar las representaciones preentrenadas: la exactitud mejora respecto a la sonda lineal.
  \item Al aumentar la tasa, los primeros pasos, con gradientes grandes procedentes de una cabeza aleatoria, desplazan mucho los parámetros del cuerpo y destruyen las representaciones aprendidas en la tarea de origen.
  \item Con un conjunto de destino pequeño, el modelo pierde entonces la ventaja de la transferencia: se comporta como una red entrenada desde una inicialización arbitraria con pocos datos, y la exactitud baja. Con tasas aún mayores, el entrenamiento puede volverse inestable.
\end{enumerate}
\resultado \textbf{la exactitud sube al principio y después cae}, porque una tasa grande anula la regularización implícita que supone permanecer cerca de $\theta_0$. Es una forma de olvido catastrófico (\cref{tema:3}).
\comprobacion la receta de~\cite[§14.2]{d2l} usa una tasa pequeña en el cuerpo y una diez veces mayor solo en la cabeza nueva, y ULMFiT reduce aún más la tasa en las capas inferiores.
\end{solucion}

\begin{solucion}{eje:1-convexa}
\planteamiento aplicar la \cref{prop:sonda-convexa} y las definiciones del tema.
\pasos
\begin{enumerate}[label=(\alph*)]
  \item \textbf{Falsa.} La convexidad garantiza que todo mínimo local es global, pero no la unicidad. Sumar el mismo vector $w$ a todas las filas de $V$ y la misma constante a todas las componentes de $b$ suma a cada logit $w\cdot\phi+\text{cte}$, igual para todas las clases, y no cambia la softmax (módulo~1, tema~4): si existe un minimizador, existen infinitos. Además, si los datos son separables, el mínimo no se alcanza.
  \item \textbf{Falsa.} Con una capa oculta, los logits dependen de los productos de los pesos de las dos capas y de una no linealidad; la pérdida no es convexa (por ejemplo, permutar las neuronas ocultas da otro mínimo con el mismo valor, y el punto medio entre ambos no lo es en general).
  \item \textbf{Falsa.} Por la \cref{def:discriminativas}, $\eta^{l-1}=\eta^{l}/2{,}6<\eta^{l}$: la mayor es la de la última capa.
\end{enumerate}
\resultado \textbf{(a) falsa; (b) falsa; (c) falsa.}
\end{solucion}

\begin{solucion}{eje:1-discriminativas}
\datos $L=12$, $\eta^{12}=0{,}01$, factor $2{,}6$.
\planteamiento $\eta^{l}=\eta^{L}/2{,}6^{\,L-l}$.
\pasos
\begin{enumerate}[label=\arabic*.]
  \item $\eta^{11}=0{,}01/2{,}6=3{,}846\cdot 10^{-3}$.
  \item $\eta^{10}=0{,}01/2{,}6^{2}=0{,}01/6{,}76=1{,}479\cdot 10^{-3}$.
  \item $\eta^{1}=0{,}01/2{,}6^{11}=0{,}01/36\,703=2{,}725\cdot 10^{-7}$.
  \item Se busca el menor $L$ con $0{,}01/2{,}6^{L-1}<10^{-6}$, es decir, $2{,}6^{L-1}>10^{4}$, o $L-1>\ln 10^{4}/\ln 2{,}6=9{,}210/0{,}9555=9{,}639$.
  \item El menor entero es $L-1=10$, luego $L=11$.
\end{enumerate}
\resultado \textbf{$\eta^{11}=3{,}846\cdot 10^{-3}$, $\eta^{10}=1{,}479\cdot 10^{-3}$, $\eta^{1}=2{,}725\cdot 10^{-7}$; basta con $L=11$ capas.}
\comprobacion con $L=10$, $\eta^{1}=0{,}01/2{,}6^{9}=0{,}01/5430=1{,}842\cdot 10^{-6}>10^{-6}$; con $L=11$, $0{,}01/14\,117=7{,}084\cdot 10^{-7}<10^{-6}$.
\end{solucion}

\begin{solucion}{eje:1-stlr}
\datos $T=1000$, $\mathrm{cf}=0{,}1$, $\rho=32$, $\eta_{\max}=0{,}01$.
\planteamiento $c=\lfloor 1000\cdot 0{,}1\rfloor=100$; $c(1/\mathrm{cf}-1)=100\cdot 9=900$; $\eta_t=0{,}01\,(1+31p_t)/32$.
\pasos
\begin{enumerate}[label=\arabic*.]
  \item $t=0$: $p=0$, $\eta=0{,}01/32=3{,}125\cdot 10^{-4}$.
  \item $t=50$: $p=50/100=0{,}5$, $\eta=0{,}01\,(1+15{,}5)/32=0{,}01\cdot 16{,}5/32=5{,}156\cdot 10^{-3}$.
  \item $t=100$: $p=1-0/900=1$, $\eta=0{,}01\cdot 32/32=0{,}01$.
  \item $t=550$: $p=1-450/900=0{,}5$, $\eta=5{,}156\cdot 10^{-3}$.
  \item $t=1000$: $p=1-900/900=0$, $\eta=3{,}125\cdot 10^{-4}$.
\end{enumerate}
\resultado \textbf{$3{,}125\cdot 10^{-4}$; $5{,}156\cdot 10^{-3}$; $0{,}01$; $5{,}156\cdot 10^{-3}$; $3{,}125\cdot 10^{-4}$.}
\comprobacion la tasa sube en 100 iteraciones y baja en 900, con valores simétricos en $p$ ($t=50$ y $t=550$ tienen $p=0{,}5$), como predice la \cref{prop:stlr}.
\end{solucion}

\begin{solucion}{eje:1-media}
\datos $\eta(t)=\eta_{\max}(1+p(t)(\rho-1))/\rho$, con $p$ lineal a trozos; $c=T\,\mathrm{cf}$.
\planteamiento por linealidad de la integral, basta calcular el valor medio de $p$.
\pasos
\begin{enumerate}[label=\arabic*.]
  \item En $[0,c]$, $p(t)=t/c$ y $\int_0^{c}p\,dt=c/2$.
  \item En $[c,T]$, $p$ decrece linealmente de $1$ a $0$ (\cref{prop:stlr}) y $\int_c^{T}p\,dt=(T-c)/2$, el área de un triángulo de base $T-c$ y altura~1.
  \item Valor medio de $p$: $\frac1T\bigl(\frac c2+\frac{T-c}{2}\bigr)=\frac12$, independiente de $c$ y, por tanto, de $\mathrm{cf}$.
  \item Valor medio de $\eta$: $\frac{\eta_{\max}}{\rho}\bigl(1+\frac12(\rho-1)\bigr)=\eta_{\max}\frac{\rho+1}{2\rho}$.
  \item Con los valores por defecto: $0{,}01\cdot 33/64=5{,}156\cdot 10^{-3}$.
\end{enumerate}
\resultado \textbf{$\bar\eta=\eta_{\max}(\rho+1)/(2\rho)$; con los valores por defecto, $5{,}156\cdot 10^{-3}$.}
\comprobacion coincide con el valor en $p=0{,}5$ del \cref{eje:1-stlr}, como corresponde a una función afín de $p$ cuya media es $0{,}5$.
\end{solucion}

% ---------------------------------------------------------------------
\section{Tema 2}

\begin{solucion}{eje:2-base}
\planteamiento comparar lo que aprendió el modelo base (la distribución del corpus) con lo que se le pide (un turno de respuesta corto que termina).
\pasos
\begin{enumerate}[label=\arabic*.]
  \item El preentrenamiento minimiza la entropía cruzada sobre texto continuo: el modelo aprende $p(x_t\mid x_{1:t-1})$ tal como aparece en el corpus.
  \item En ese corpus casi no aparece el patrón «p: \dots\ r:» seguido de una respuesta corta y un cierre. La continuación más probable es más texto con el aspecto del corpus, no «paris».
  \item El dato «París es la capital de Francia» está en los pesos, pero asociado a frases del corpus, no a la pregunta en formato de chat.
  \item El modelo base nunca recibió un token de fin de turno como objetivo: aunque empezara bien, no sabría parar.
\end{enumerate}
\resultado \textbf{le falta el formato, no el conocimiento.} El ajuste supervisado le enseña a responder en ese formato y a cerrar el turno con el token de fin.
\comprobacion es la hipótesis de la alineación superficial (LIMA). En la práctica, 600 pasos de SFT bastan para que el mismo mini-GPT responda y emita el token de fin en el $100\,\%$ de los ejemplos de prueba.
\end{solucion}

\begin{solucion}{eje:2-ignore}
\planteamiento recordar cómo calcula PyTorch la media de la entropía cruzada con \texttt{ignore\_\allowbreak index}.
\pasos
\begin{enumerate}[label=\arabic*.]
  \item Con reducción media, PyTorch calcula la suma de las pérdidas de las posiciones válidas dividida por el número de posiciones válidas.
  \item Una posición con etiqueta $-100$ no suma en el numerador ni cuenta en el denominador, y su gradiente respecto a los logits es nulo (\cref{prop:grad-mascara} con $m_t=0$).
  \item El relleno solo sirve para igualar longitudes dentro de un lote; esos tokens no forman parte de ningún ejemplo.
  \item Sin $-100$, el modelo aprendería a predecir relleno y la media se diluiría con posiciones irrelevantes.
\end{enumerate}
\resultado \textbf{la posición se ignora por completo (numerador, denominador y gradiente)}; en el relleno se usa porque esos tokens no son datos.
\comprobacion si todas las etiquetas de la instrucción y del relleno valen $-100$, la pérdida coincide con la \cref{def:sft}.
\end{solucion}

\begin{solucion}{eje:2-mascara}
\datos instrucción $(2{,}1;\,1{,}8;\,0{,}9;\,1{,}1)$; respuesta $(0{,}6;\,0{,}2;\,0{,}4)$; $T=7$, $\sum_t m_t=3$.
\planteamiento sin máscara se divide la suma de las 7 pérdidas entre 7; con máscara, la suma de las de la respuesta entre~3.
\pasos
\begin{enumerate}[label=\arabic*.]
  \item Suma de la instrucción: $2{,}1+1{,}8+0{,}9+1{,}1=5{,}9$.
  \item Suma de la respuesta: $0{,}6+0{,}2+0{,}4=1{,}2$.
  \item Sin máscara: $(5{,}9+1{,}2)/7=7{,}1/7=1{,}014$.
  \item Con máscara: $1{,}2/3=0{,}400$.
\end{enumerate}
\resultado \textbf{sin máscara $1{,}014$ nats; con máscara $0{,}400$ nats.}
\comprobacion $1{,}014$ está entre la media de la instrucción ($5{,}9/4=1{,}475$) y la de la respuesta ($0{,}400$), más cerca de la primera porque tiene más tokens.
\end{solucion}

\begin{solucion}{eje:2-tokenizar}
\datos ejemplo~1: $[5,9,2]$ y $[7,3,0]$; ejemplo~2: $[4,8]$ y $[6,0]$; relleno $=1$.
\planteamiento la etiqueta de la posición $t$ es el token $t+1$; la máscara vale~1 si esa etiqueta pertenece a la respuesta.
\pasos
\begin{enumerate}[label=\arabic*.]
  \item Concatenar: $[5,9,2,7,3,0]$ (6 tokens) y $[4,8,6,0]$ (4 tokens). Longitud máxima: 6.
  \item Rellenar: $[4,8,6,0,1,1]$.
  \item \texttt{input\_\allowbreak ids} (sin el último token): $[5,9,2,7,3]$ y $[4,8,6,0,1]$.
  \item \texttt{labels} (sin el primer token): $[9,2,7,3,0]$ y $[8,6,0,1,1]$.
  \item Máscara. Ejemplo~1: $9$ y $2$ son de la instrucción; $7$, $3$ y $0$, de la respuesta: $[0,0,1,1,1]$. Ejemplo~2: $8$ es de la instrucción; $6$ y $0$, de la respuesta; los dos $1$ son relleno: $[0,1,1,0,0]$.
\end{enumerate}
\resultado \textbf{máscaras $[0,0,1,1,1]$ y $[0,1,1,0,0]$: puntúan 3 y 2 posiciones.}
\comprobacion cada ejemplo puntúa tantas posiciones como tokens tiene su respuesta, incluido el de fin.
\end{solucion}

\begin{solucion}{eje:2-microlote}
\datos tamaño~2, $c=1$, 2 pasos de acumulación y
\begin{gather*}
\ln p^{(1)}=(-2{,}0;\,-1{,}5;\,-0{,}5;\,-1{,}0;\,-0{,}1),\quad m^{(1)}=(0,0,1,1,1);\\
\ln p^{(2)}=(-1{,}2;\,-0{,}4;\,-0{,}2;\,-3{,}0;\,-3{,}0),\quad m^{(2)}=(0,1,1,0,0).
\end{gather*}
\planteamiento sumar $-\ln p$ donde $m=1$ y dividir por tamaño, $c$ y pasos de acumulación (\cref{prop:acumulacion}).
\pasos
\begin{enumerate}[label=\arabic*.]
  \item Secuencia~1: $-(-0{,}5-1{,}0-0{,}1)=1{,}6$.
  \item Secuencia~2: $-(-0{,}4-0{,}2)=0{,}6$.
  \item Suma enmascarada: $1{,}6+0{,}6=2{,}2$.
  \item Entre el tamaño del microlote y $c$: $2{,}2/2/1=1{,}1$.
  \item Entre los pasos de acumulación: $1{,}1/2=0{,}55$.
  \item $J_{\mathrm{tok}}=2{,}2/(3+2)=0{,}44$.
\end{enumerate}
\resultado \textbf{pérdida del microlote $0{,}55$; media por token $0{,}44$.}
\comprobacion con dos microlotes iguales, $0{,}55+0{,}55=1{,}1$, la pérdida por secuencia sin acumulación. Las dos normalizaciones difieren porque las respuestas tienen longitudes distintas (3 y 2).
\end{solucion}

\begin{solucion}{eje:2-empaquetar}
\datos ejemplos de 7, 5 y 9 tokens; un fin tras cada uno; trozos de 5; alternativa con relleno a 10.
\planteamiento contar los tokens totales, dividir entre la longitud del trozo y comparar posiciones útiles con posiciones totales.
\pasos
\begin{enumerate}[label=\arabic*.]
  \item Tokens con fin: $8+6+10=24$.
  \item Trozos completos: $24=4\cdot 5+4$. Salen 4 secuencias y se pierden 4 tokens.
  \item Empaquetando se aprovechan $20/24=5/6=0{,}8333$ de los tokens, y todas las posiciones de cada secuencia son reales.
  \item Con relleno: $3\cdot 10=30$ posiciones, de las que 24 son reales: $24/30=0{,}80$.
\end{enumerate}
\resultado \textbf{4 secuencias y 4 tokens perdidos ($5/6\approx 83\,\%$ aprovechado); con relleno, solo el $80\,\%$ de las posiciones son útiles.}
\comprobacion $4\cdot 5+4=24$ y $24+6=30$.
\end{solucion}

\begin{solucion}{eje:2-multiturno}
\datos secuencia $u^{(1)},y^{(1)},\dots,u^{(K)},y^{(K)}$; máscara 1 exactamente en los tokens de las $y^{(k)}$.
\planteamiento aplicar la regla de la cadena a la secuencia completa y agrupar los factores por segmentos, como en la \cref{prop:sft-condicional}.
\pasos
\begin{enumerate}[label=\arabic*.]
  \item El denominador es $\sum_tm_t=\sum_km_k$.
  \item El numerador es $\sum_{t:m_t=1}\ln p_\theta(x_t\mid x_{1:t-1})$. Las posiciones con $m_t=1$ son las de las respuestas, y para el token $j$ de $y^{(k)}$ el contexto $x_{1:t-1}$ es $(h^{(k)},y^{(k)}_{1:j-1})$.
  \item Agrupando por respuestas: $\sum_k\sum_{j=1}^{m_k}\ln p_\theta(y^{(k)}_j\mid h^{(k)},y^{(k)}_{1:j-1})$.
  \item Por la regla de la cadena condicionada a $h^{(k)}$, cada suma interior es $\ln p_\theta(y^{(k)}\mid h^{(k)})$.
  \item Sustituyendo y cambiando el signo se obtiene la fórmula del enunciado.
\end{enumerate}
\resultado \textbf{$\Lsft=-\frac{1}{\sum_km_k}\sum_k\ln p_\theta(y^{(k)}\mid h^{(k)})$.} El ajuste maximiza la verosimilitud de cada respuesta del asistente dado todo el historial anterior, incluidas sus propias respuestas previas; los turnos del usuario solo actúan como contexto.
\comprobacion con $K=1$ se recupera la \cref{prop:sft-condicional}, con $h^{(1)}=u^{(1)}$.
\end{solucion}

% ---------------------------------------------------------------------
\section{Tema 3}

\begin{solucion}{eje:3-remedios}
\planteamiento repasar los remedios de la \cref{tema:3} y de los temas~1, 4 y~5.
\pasos
\begin{enumerate}[label=\arabic*.]
  \item Sobre los datos: repetición de datos generales (\cref{def:replay}).
  \item Sobre la pérdida: penalizar el alejamiento de $\theta_0$ (muelle, EWC) o de las salidas del modelo de referencia (divergencia KL, módulo~21).
  \item Sobre el optimizador: tasa pequeña, pocas épocas, tasas discriminativas y descongelación gradual (\cref{tema:1}).
  \item Sobre los parámetros entrenados: métodos PEFT y LoRA, que dejan $\theta_0$ congelado.
\end{enumerate}
\resultado \textbf{solo los métodos que congelan $\theta_0$ (PEFT, LoRA) permiten recuperar exactamente el modelo base}: basta con desactivar el adaptador, porque los pesos originales no se han modificado.
\comprobacion en el \cref{ej:olvido-practica}, desactivar el adaptador LoRA devuelve la pérdida de texto exacta del modelo base, $0{,}192$ nats.
\end{solucion}

\begin{solucion}{eje:3-muelle}
\datos $\theta_0=1$; tarea $(\theta-3)^{2}$; muelle $\lambda(\theta-1)^{2}$; $\lambda=2$.
\planteamiento derivar e igualar a cero; la pérdida es una parábola abierta hacia arriba.
\pasos
\begin{enumerate}[label=\arabic*.]
  \item $\mathcal{L}'(\theta)=2(\theta-3)+2\lambda(\theta-1)$.
  \item $\mathcal{L}'=0\iff\theta-3+\lambda\theta-\lambda=0\iff\theta(1+\lambda)=3+\lambda\iff\theta^{\ast}=(3+\lambda)/(1+\lambda)$.
  \item $\mathcal{L}''=2+2\lambda>0$: es un mínimo.
  \item $\lambda=2$: $\theta^{\ast}=5/3\approx 1{,}667$.
  \item Olvido: $(5/3-1)^{2}=(2/3)^{2}=4/9\approx 0{,}4444$.
  \item Error en la tarea: $(5/3-3)^{2}=(-4/3)^{2}=16/9\approx 1{,}778$.
  \item Límite: $\theta^{\ast}=(3/\lambda+1)/(1/\lambda+1)\to 1$ cuando $\lambda\to\infty$.
\end{enumerate}
\resultado \textbf{$\theta^{\ast}=(3+\lambda)/(1+\lambda)$; con $\lambda=2$, $\theta^{\ast}=5/3$, olvido $4/9$ y error $16/9$. Si $\lambda\to\infty$, $\theta^{\ast}\to 1$: no olvida, pero tampoco aprende.}
\comprobacion por el \cref{cor:frontera}, $\sqrt{4/9}+\sqrt{16/9}=2/3+4/3=2=\abs{3-1}$.
\end{solucion}

\begin{solucion}{eje:3-ewc}
\datos $\theta_A=(1,1)$; $F=(4;\,0{,}25)$; $\lambda=2$; $\mathcal{L}_B=(\theta_1-3)^{2}+(\theta_2-3)^{2}$.
\planteamiento la pérdida se separa por coordenadas: $(\theta_i-3)^{2}+\frac{\lambda F_i}{2}(\theta_i-1)^{2}$. Por la \cref{prop:muelle} con muelle $\frac{\lambda F_i}{2}$, $\theta_i^{\ast}=(3+\frac{\lambda F_i}{2})/(1+\frac{\lambda F_i}{2})$.
\pasos
\begin{enumerate}[label=\arabic*.]
  \item Coordenada~1: $\lambda F_1/2=2\cdot 4/2=4$; $\theta^{\ast}_1=(3+4)/(1+4)=7/5=1{,}4$.
  \item Coordenada~2: $\lambda F_2/2=2\cdot 0{,}25/2=0{,}25$; $\theta^{\ast}_2=3{,}25/1{,}25=2{,}6$.
  \item Penalización: $\frac12\bigl(4\cdot 0{,}4^{2}+0{,}25\cdot 1{,}6^{2}\bigr)=\frac12(0{,}64+0{,}64)=0{,}64$.
  \item Error en la tarea nueva: $(1{,}4-3)^{2}+(2{,}6-3)^{2}=2{,}56+0{,}16=2{,}72$.
\end{enumerate}
\resultado \textbf{$\theta^{\ast}=(1{,}4;\,2{,}6)$; penalización $0{,}64$; $\mathcal{L}_B(\theta^{\ast})=2{,}72$.} Se mueve mucho más $\theta_2$, cuya información de Fisher es 16 veces menor: la tarea antigua apenas depende de él.
\comprobacion derivada en la coordenada~1: $2(1{,}4-3)+2\cdot 4\cdot(1{,}4-1)=-3{,}2+3{,}2=0$; en la~2: $2(2{,}6-3)+2\cdot 0{,}25\cdot(2{,}6-1)=-0{,}8+0{,}8=0$.
\end{solucion}

\begin{solucion}{eje:3-replay}
\datos $\mathcal{L}_B=(\theta-3)^{2}$, $\mathcal{L}_A=(\theta-1)^{2}$, $\beta=0{,}25$.
\planteamiento derivar la mezcla e igualar a cero.
\pasos
\begin{enumerate}[label=\arabic*.]
  \item $\frac{d}{d\theta}\bigl[(1-\beta)(\theta-3)^{2}+\beta(\theta-1)^{2}\bigr]=2(1-\beta)(\theta-3)+2\beta(\theta-1)$.
  \item Igualando a cero: $\theta=(1-\beta)\cdot 3+\beta\cdot 1=3-2\beta$.
  \item $\beta=0{,}25$: $\theta^{\ast}=3-0{,}5=2{,}5$.
  \item Con $\lambda=\beta/(1-\beta)$, la \cref{prop:muelle} da $\frac{3+\lambda}{1+\lambda}=\frac{3(1-\beta)+\beta}{(1-\beta)+\beta}=3-2\beta$, tras multiplicar numerador y denominador por $1-\beta$.
  \item Para $\beta=0{,}25$, $\lambda=1/3$ y $(3+1/3)/(1+1/3)=(10/3)/(4/3)=2{,}5$.
\end{enumerate}
\resultado \textbf{$\theta^{\ast}(\beta)=3-2\beta$; con $\beta=0{,}25$, $\theta^{\ast}=2{,}5$, igual que el muelle con $\lambda=1/3$.}
\comprobacion la mezcla es $(1-\beta)\bigl[\mathcal{L}_B+\tfrac{\beta}{1-\beta}\mathcal{L}_A\bigr]$, y el factor positivo $1-\beta$ no cambia el minimizador.
\end{solucion}

\begin{solucion}{eje:3-fisher}
\datos $p_z(1)=\sigma(z)$, $p_z(0)=1-\sigma(z)$; $\sigma'(z)=\sigma(z)(1-\sigma(z))$.
\planteamiento calcular la puntuación y su segundo momento (\cref{def:fisher}).
\pasos
\begin{enumerate}[label=\arabic*.]
  \item $\ln p_z(x)=x\ln\sigma(z)+(1-x)\ln(1-\sigma(z))$ para $x\in\{0,1\}$.
  \item Puntuación: $\frac{d}{dz}\ln p_z(x)=x\frac{\sigma'}{\sigma}-(1-x)\frac{\sigma'}{1-\sigma}=x(1-\sigma)-(1-x)\sigma=x-\sigma(z)$.
  \item $\Fis(z)=\E[(x-\sigma)^{2}]=\operatorname{Var}(x)=\sigma(z)(1-\sigma(z))$, porque $\E[x]=\sigma(z)$.
  \item $z=0$: $\sigma=0{,}5$ y $\Fis=0{,}25$.
  \item $z=\ln 9$: $\sigma=1/(1+1/9)=0{,}9$ y $\Fis=0{,}9\cdot 0{,}1=0{,}09$.
\end{enumerate}
\resultado \textbf{$\Fis(z)=\sigma(z)(1-\sigma(z))$; vale $0{,}25$ en $z=0$ y $0{,}09$ en $z=\ln 9$.} El parámetro con $z=0$ quedaría unas $2{,}8$ veces más rígido: una predicción incierta es más sensible a cambios del logit.
\comprobacion por la \cref{prop:fisher-hess}, $-\frac{d^{2}}{dz^{2}}\ln p_z(x)=\frac{d}{dz}\sigma(z)=\sigma(1-\sigma)$, que no depende de $x$ y coincide con el resultado.
\end{solucion}

\begin{solucion}{eje:3-gbc}
\datos $\mathcal{L}_B=\frac12(\theta-\theta_B)\T H(\theta-\theta_B)$, $H=Q\Lambda Q\T$; penalización $\frac{\lambda}{2}\norm{\theta-\theta_0}^{2}$.
\planteamiento anular el gradiente y diagonalizar con la base de autovectores.
\pasos
\begin{enumerate}[label=\arabic*.]
  \item (a) El gradiente es $H(\theta-\theta_B)+\lambda(\theta-\theta_0)$. Se anula si $(H+\lambda I)\theta=H\theta_B+\lambda\theta_0$.
  \item $H+\lambda I$ tiene autovalores $h_i+\lambda>0$, luego es definida positiva e invertible, y la función es estrictamente convexa: $\theta^{\ast}=(H+\lambda I)^{-1}(H\theta_B+\lambda\theta_0)$ es el único minimizador.
  \item (b) Restando $\theta_0=(H+\lambda I)^{-1}(H+\lambda I)\theta_0$: $\theta^{\ast}-\theta_0=(H+\lambda I)^{-1}H(\theta_B-\theta_0)$.
  \item Con $H=Q\Lambda Q\T$, $(H+\lambda I)^{-1}H=Q(\Lambda+\lambda I)^{-1}\Lambda Q\T$. Multiplicando por $q_i\T$: $q_i\T(\theta^{\ast}-\theta_0)=\frac{h_i}{h_i+\lambda}\,q_i\T(\theta_B-\theta_0)$.
  \item (c) $H$ ya es diagonal, $\theta_0=0$: $\theta^{\ast}_1=\frac{10}{11}\cdot 1=0{,}9091$ y $\theta^{\ast}_2=\frac{0{,}1}{1{,}1}\cdot 1=0{,}09091$.
\end{enumerate}
\resultado \textbf{$\theta^{\ast}=(H+\lambda I)^{-1}(H\theta_B+\lambda\theta_0)$; factor $h_i/(h_i+\lambda)$ por dirección; en (c), $\theta^{\ast}=(0{,}9091;\,0{,}09091)$.} En la dirección de curvatura alta, importante para la tarea nueva, el modelo casi alcanza $\theta_B$; en la de curvatura baja, la penalización domina y el parámetro apenas se aleja de $\theta_0$.
\comprobacion con $\lambda=0$ el factor es~1 ($\theta^{\ast}=\theta_B$); con $\lambda\to\infty$ tiende a~0 ($\theta^{\ast}\to\theta_0$), como en el \cref{cor:frontera}.
\end{solucion}

% ---------------------------------------------------------------------
\section{Tema 4}

\begin{solucion}{eje:4-congelar}
\datos $h=768$; $K=2$ clases con sesgo; cuerpo de $1{,}1\cdot 10^{8}$ parámetros.
\planteamiento la cabeza tiene $Kh+K$ parámetros (\cref{def:sonda}).
\pasos
\begin{enumerate}[label=\arabic*.]
  \item (a) $2\cdot 768+2=1538$ parámetros; fracción $1538/(1{,}1\cdot 10^{8}+1538)=1{,}398\cdot 10^{-5}$, un $0{,}0014\,\%$.
  \item (b) Es un método selectivo: se entrena un subconjunto de los parámetros (solo la cabeza). Es la extracción de características o sonda lineal.
  \item (c) En general, menos exactitud: las representaciones no se adaptan a la tarea. La diferencia es pequeña si la tarea de destino se parece a la de origen, de modo que las representaciones ya separan bien las clases, o si hay muy pocos datos, porque entonces el ajuste completo sobreajusta.
\end{enumerate}
\resultado \textbf{1538 parámetros ($0{,}0014\,\%$); método selectivo; exactitud en general menor, con diferencia pequeña si las tareas se parecen o hay pocos datos.}
\comprobacion por la \cref{prop:sonda-convexa}, este problema es convexo, lo que lo hace estable y barato aunque menos expresivo.
\end{solucion}

\begin{solucion}{eje:4-adaptadores}
\datos $L=12$; $d=768$; $m=64$; 2 adaptadores por capa; modelo de $1{,}1\cdot 10^{8}$ parámetros; 20 tareas.
\planteamiento por la \cref{prop:adaptador}, cada adaptador tiene $2md+m+d$ parámetros.
\pasos
\begin{enumerate}[label=\arabic*.]
  \item Por adaptador: $2\cdot 64\cdot 768+64+768=98\,304+832=99\,136$.
  \item Total: $2\cdot 12\cdot 99\,136=2\,379\,264$.
  \item Porcentaje: $2\,379\,264/1{,}1\cdot 10^{8}=0{,}02163$, un $2{,}16\,\%$.
  \item (b) Con adaptadores: un modelo base y 20 juegos de adaptadores, $1{,}1\cdot 10^{8}+20\cdot 2\,379\,264=1{,}1\cdot 10^{8}+47\,585\,280=1{,}576\cdot 10^{8}$ parámetros.
  \item Con copias completas: $20\cdot 1{,}1\cdot 10^{8}=2{,}2\cdot 10^{9}$ parámetros.
\end{enumerate}
\resultado \textbf{$2\,379\,264$ parámetros ($2{,}16\,\%$); para 20 tareas, $1{,}58\cdot 10^{8}$ frente a $2{,}2\cdot 10^{9}$ parámetros, unas 14 veces menos.}
\comprobacion $2{,}2\cdot 10^{9}/1{,}576\cdot 10^{8}=13{,}96$. El porcentaje es del orden del $3{,}6\,\%$ que citan Houlsby \emph{et al.}, que usan otros tamaños de cuello.
\end{solucion}

\begin{solucion}{eje:4-prefijo}
\datos $L=24$; $d=1024$; $p=20$.
\planteamiento \cref{def:prompt-tuning} ($pd$) y \cref{def:prefix} ($2Lpd$).
\pasos
\begin{enumerate}[label=\arabic*.]
  \item Ajuste de \emph{prompt}: $20\cdot 1024=20\,480$.
  \item Ajuste de prefijo: $2\cdot 24\cdot 20\cdot 1024=983\,040$.
  \item Cociente: $983\,040/20\,480=48$.
\end{enumerate}
\resultado \textbf{$20\,480$ frente a $983\,040$ parámetros; el prefijo tiene $48=2L$ veces más.}
\comprobacion el cociente es $2Lpd/(pd)=2L$: un par de claves y valores por capa frente a un único \emph{embedding} de entrada.
\end{solucion}

\begin{solucion}{eje:4-subespacio}
\datos $\theta_0=(1,0,2)$; $P=(1,2,-1)\T$; $\theta^{\star}=(2,3,1)$; $\mathcal{L}=\norm{\theta-\theta^{\star}}^{2}$.
\planteamiento $\mathcal{L}(z)=\norm{\theta_0-\theta^{\star}+zP}^{2}$ es una parábola en $z$.
\pasos
\begin{enumerate}[label=\arabic*.]
  \item $\theta_0-\theta^{\star}=(-1,-3,1)$; $\norm{P}^{2}=1+4+1=6$; $P\T(\theta_0-\theta^{\star})=-1-6-1=-8$; $\norm{\theta_0-\theta^{\star}}^{2}=1+9+1=11$.
  \item $\mathcal{L}(z)=6z^{2}-16z+11$. Mínimo en $z^{\ast}=16/12=4/3$.
  \item $\theta=\theta_0+\frac43P=(1+\frac43,\ \frac83,\ 2-\frac43)=(7/3,\ 8/3,\ 2/3)$.
  \item (b) $\mathcal{L}(4/3)=6\cdot\frac{16}{9}-\frac{64}{3}+11=\frac{32}{3}-\frac{64}{3}+\frac{33}{3}=\frac13$. El ajuste completo alcanza $\theta=\theta^{\star}$ y pérdida~0.
  \item (c) En $z=0$, $\nabla_\theta\mathcal{L}=2(\theta_0-\theta^{\star})=(-2,-6,2)$ y $P\T\nabla_\theta\mathcal{L}=-2-12-2=-16$, que coincide con $\mathcal{L}'(0)=-16$.
\end{enumerate}
\resultado \textbf{$z^{\ast}=4/3$, $\theta=(7/3;\,8/3;\,2/3)$; pérdida residual $1/3$ frente a~0; el gradiente vale $-16$ por los dos caminos.}
\comprobacion el residuo $\theta^{\star}-\theta=(-1/3,1/3,1/3)$ es ortogonal a $P$: $-\frac13+\frac23-\frac13=0$, como debe ocurrir en una proyección ortogonal.
\end{solucion}

\begin{solucion}{eje:4-proyeccion}
\datos $\theta=\theta_0+Pz$, $P\in\R^{D\times d'}$.
\planteamiento usar la definición de espacio columna y la \cref{prop:grad-subespacio}.
\pasos
\begin{enumerate}[label=\arabic*.]
  \item (a) $\{Pz:z\in\R^{d'}\}$ es, por definición, $\Col(P)$, un subespacio vectorial de dimensión $\rango P$. Trasladado por $\theta_0$ es el subespacio afín $\theta_0+\Col(P)$.
  \item (b) Un paso en $z$ es $z\gets z-\eta\nabla_z\mathcal{L}=z-\eta P\T\nabla_\theta\mathcal{L}$. Multiplicando por $P$, $\theta$ cambia en $P\cdot(-\eta P\T\nabla_\theta\mathcal{L})=-\eta PP\T\nabla_\theta\mathcal{L}$.
  \item (c) Si $P\T P=I$, la matriz $\Pi=PP\T$ es simétrica e idempotente ($\Pi^{2}=P(P\T P)P\T=\Pi$), con imagen $\Col(P)$: es el proyector ortogonal sobre $\Col(P)$.
  \item El paso del ajuste completo es $-\eta\nabla_\theta\mathcal{L}$; su proyección es $\Pi(-\eta\nabla_\theta\mathcal{L})=-\eta PP\T\nabla_\theta\mathcal{L}$, que coincide con (b).
\end{enumerate}
\resultado \textbf{los modelos alcanzables forman $\theta_0+\Col(P)$; el paso en el subespacio es $-\eta PP\T\nabla_\theta\mathcal{L}$, la proyección ortogonal del paso completo cuando $P$ tiene columnas ortonormales.}
\comprobacion en el \cref{eje:4-subespacio}, $P/\sqrt6$ es unitario y $\Pi\nabla_\theta\mathcal{L}=\frac{-16}{6}P$, proporcional a $P$, como corresponde a una proyección sobre una recta.
\end{solucion}

% ---------------------------------------------------------------------
\section{Tema 5}

\begin{solucion}{eje:5-vf}
\planteamiento aplicar la \cref{def:lora}, el \cref{cor:init} y la \cref{prop:fusion}.
\pasos
\begin{enumerate}[label=(\alph*)]
  \item \textbf{Falsa.} $W$ no recibe actualizaciones: el optimizador solo contiene $A$ y $B$. $W'$ cambia porque cambia $BA$. Por eso desactivar el adaptador devuelve exactamente el modelo base (\cref{ej:olvido-practica}).
  \item \textbf{Falsa.} Con $A=B=0$, $\partial\mathcal{L}/\partial B=s\,G(Ax)\T=0$ y $\partial\mathcal{L}/\partial A=s\,B\T Gx\T=0$: ningún parámetro se mueve (\cref{eje:5-init}).
  \item \textbf{Verdadera.} $W'$ tiene la misma forma que $W$ y $W'x$ cuesta las mismas operaciones (\cref{prop:fusion}).
\end{enumerate}
\resultado \textbf{(a) falsa; (b) falsa; (c) verdadera.}
\comprobacion (c) solo vale tras fusionar; sin fusionar, cada capa calcula además $B(Ax)$, con $2r(d+k)$ operaciones más.
\end{solucion}

\begin{solucion}{eje:5-contar}
\datos (a) $d=1024$, $k=4096$, $r=8$. (b) 12 capas, 4 matrices de $768\times 768$ por capa, $r=8$.
\planteamiento \cref{prop:lora-param}: $dk$ frente a $r(d+k)$.
\pasos
\begin{enumerate}[label=\arabic*.]
  \item (a) Completo: $1024\cdot 4096=4\,194\,304$.
  \item LoRA: $8\cdot(1024+4096)=8\cdot 5120=40\,960$.
  \item Porcentaje: $40\,960/4\,194\,304=9{,}766\cdot 10^{-3}$, un $0{,}977\,\%$.
  \item (b) LoRA por matriz: $8\cdot 1536=12\,288$. Total: $12\cdot 4\cdot 12\,288=589\,824$.
  \item Matrices completas: $48\cdot 768^{2}=48\cdot 589\,824=28\,311\,552$.
  \item Porcentaje: $589\,824/28\,311\,552=0{,}02083$, un $2{,}08\,\%$.
\end{enumerate}
\resultado \textbf{(a) $4\,194\,304$ parámetros; LoRA entrena $40\,960$ ($0{,}98\,\%$). (b) $589\,824$ frente a $28\,311\,552$: $2{,}08\,\%$.}
\comprobacion (a) $r(1/d+1/k)=8(1/1024+1/4096)=8\cdot 1{,}221\cdot 10^{-3}=9{,}766\cdot 10^{-3}$. (b) con $d=k$ la fracción es $2r/d=16/768=0{,}02083$; todo LoRA ocupa lo mismo que una sola matriz de $768\times 768$.
\end{solucion}

\begin{solucion}{eje:5-fusion}
\datos $W=I_2$, $B=(1,0)\T$, $A=(2,1)$, $r=1$, $\alpha=2$, $x=(1,1)$.
\planteamiento $W'=W+s\,BA$ con $s=\alpha/r$, y comprobación con $h=Wx+s\,B(Ax)$.
\pasos
\begin{enumerate}[label=\arabic*.]
  \item $s=2/1=2$.
  \item $BA=\begin{pmatrix}1\cdot 2&1\cdot 1\\0\cdot 2&0\cdot 1\end{pmatrix}=\begin{pmatrix}2&1\\0&0\end{pmatrix}$; $2BA=\begin{pmatrix}4&2\\0&0\end{pmatrix}$.
  \item $W'=I+2BA=\begin{pmatrix}5&2\\0&1\end{pmatrix}$.
  \item $W'x=(5+2;\,0+1)=(7,1)$.
\end{enumerate}
\resultado \textbf{$W'=\begin{pmatrix}5&2\\0&1\end{pmatrix}$ y $W'x=(7,1)$.}
\comprobacion sin fusionar: $Ax=2+1=3$; $B(Ax)=(3,0)$; $h=(1,1)+2\cdot(3,0)=(7,1)$.
\end{solucion}

\begin{solucion}{eje:5-grad}
\datos $s=2$, $B=(0,0)\T$, $A=(1,1)$, $G=(1,-1)$, $x=(2,1)$.
\planteamiento \cref{prop:grad-lora}.
\pasos
\begin{enumerate}[label=\arabic*.]
  \item $Ax=1\cdot 2+1\cdot 1=3$ (escalar, porque $r=1$).
  \item $\partial\mathcal{L}/\partial B=2\cdot(1,-1)\T\cdot 3=(6,-6)\T$.
  \item $B\T G=0\cdot 1+0\cdot(-1)=0$, luego $\partial\mathcal{L}/\partial A=2\cdot 0\cdot x\T=(0,0)$.
\end{enumerate}
\resultado \textbf{$\partial\mathcal{L}/\partial B=(6,-6)\T$ y $\partial\mathcal{L}/\partial A=(0,0)$}: en el primer paso solo se mueve $B$.
\comprobacion dimensiones: $\partial\mathcal{L}/\partial B$ es $d\times r=2\times 1$ y $\partial\mathcal{L}/\partial A$, $r\times k=1\times 2$.
\end{solucion}

\begin{solucion}{eje:5-svd}
\datos $\Delta W=\begin{pmatrix}2&1&1\\1&2&-1\end{pmatrix}$.
\planteamiento los autovalores de $\Delta W\Delta W\T$ son $\sigma_i^{2}$, sus autovectores los $u_i$, y $v_i=\Delta W\T u_i/\sigma_i$; el \cref{teo:ey} da $\sigma_1u_1v_1\T$.
\pasos
\begin{enumerate}[label=\arabic*.]
  \item $\Delta W\Delta W\T=\begin{pmatrix}4+1+1&2+2-1\\2+2-1&1+4+1\end{pmatrix}=\begin{pmatrix}6&3\\3&6\end{pmatrix}$.
  \item $(6-\lambda)^{2}-9=0\Rightarrow\lambda_1=9$, $\lambda_2=3$; $\sigma_1=3$, $\sigma_2=\sqrt3=1{,}732$.
  \item Para $\lambda_1=9$: $-3u_1+3u_2=0\Rightarrow u_1=(1,1)/\sqrt2$.
  \item $v_1=\Delta W\T u_1/3=(3,3,0)/(3\sqrt2)=(1,1,0)/\sqrt2$.
  \item $\Delta W_1=3\cdot\frac12\begin{pmatrix}1&1&0\\1&1&0\end{pmatrix}=\begin{pmatrix}1{,}5&1{,}5&0\\1{,}5&1{,}5&0\end{pmatrix}$.
  \item Error: $\norm{\Delta W-\Delta W_1}_F=\sigma_2=\sqrt3\approx 1{,}732$. Energía conservada: $9/(9+3)=0{,}75$.
  \item Como producto: $\sigma_1u_1v_1\T=1{,}5\,(1,1)\T(1,1,0)$; se toma $B=(1{,}5;\,1{,}5)\T$ y $A=(1,1,0)$.
\end{enumerate}
\resultado \textbf{$\sigma=(3;\,\sqrt3)$; $\Delta W_1=\begin{pmatrix}1{,}5&1{,}5&0\\1{,}5&1{,}5&0\end{pmatrix}$; error $\sqrt3$; energía $75\,\%$; $B=(1{,}5;\,1{,}5)\T$, $A=(1,1,0)$.}
\comprobacion $\Delta W-\Delta W_1=\begin{pmatrix}0{,}5&-0{,}5&1\\-0{,}5&0{,}5&-1\end{pmatrix}$, de norma al cuadrado $0{,}25\cdot 4+1+1=3$; y $\norm{\Delta W}_F^{2}=12=\sigma_1^{2}+\sigma_2^{2}$.
\end{solucion}

\begin{solucion}{eje:5-rango}
\planteamiento probar que la imagen de $BA$ está contenida en la de $B$ y buscar una matriz de rango~2.
\pasos
\begin{enumerate}[label=\arabic*.]
  \item Para todo $x\in\R^{k}$, $(BA)x=B(Ax)$ por asociatividad. Con $z=Ax\in\R^{r}$, $Bz=z_1b_1+\dots+z_rb_r$, donde $b_j$ son las columnas de $B$.
  \item Luego $\img(BA)\subseteq\Col(B)$, generado por $r$ vectores: $\rango(BA)\le\rango B\le r$.
  \item Si $r<\min(d,k)$, existen matrices $d\times k$ de rango $\min(d,k)>r$ (la que tiene unos en la diagonal principal y ceros en el resto), que ningún $BA$ puede igualar.
  \item Ejemplo: $d=k=2$, $r=1$, $\Delta W=I_2$. Con $B=(b_1,b_2)\T$ y $A=(a_1,a_2)$, $BA=\begin{pmatrix}b_1a_1&b_1a_2\\b_2a_1&b_2a_2\end{pmatrix}$ y $\det(BA)=b_1a_1b_2a_2-b_1a_2b_2a_1=0$, mientras que $\det I_2=1$.
\end{enumerate}
\resultado \textbf{$\rango(BA)\le r$; $I_2$ no es representable con $r=1$.}
\comprobacion el determinante nulo es una prueba independiente de que toda $BA$ de $2\times 2$ con $r=1$ tiene rango a lo sumo~1.
\end{solucion}

\begin{solucion}{eje:5-init}
\planteamiento sustituir los valores iniciales en la \cref{prop:grad-lora} y razonar por inducción sobre los pasos.
\pasos
\begin{enumerate}[label=\arabic*.]
  \item Con $A=0$: $Ax=0$ y $\partial\mathcal{L}/\partial B=s\,G\cdot 0\T=0$. Con $B=0$: $B\T G=0$ y $\partial\mathcal{L}/\partial A=0$. Vale para todo $x$ y todo $G$.
  \item Actualización: $A\gets 0-\eta\cdot 0=0$ y $B\gets 0$. Por inducción sobre el número de pasos, $A=B=0$ siempre.
  \item Con Adam: si todos los gradientes son nulos, los momentos lo son y la actualización es nula (con $\varepsilon>0$ en el denominador no hay división por cero).
  \item Con $B=0$ y $A\ne 0$: $\partial\mathcal{L}/\partial B=s\,G(Ax)\T$, producto exterior de $G\in\R^{d}$ por $Ax\in\R^{r}$ multiplicado por $s\ne 0$.
  \item Si $G\ne 0$ y $Ax\ne 0$, existen $i$ con $G_i\ne 0$ y $j$ con $(Ax)_j\ne 0$, y la entrada $(i,j)$ vale $s\,G_i(Ax)_j\ne 0$.
  \item Así $B$ cambia en el primer paso y, desde el segundo, $B\ne 0$ y $A$ también recibe gradiente.
\end{enumerate}
\resultado \textbf{$(0,0)$ es un punto fijo del descenso de gradiente; con $B=0$ y $A\ne 0$, $B$ se mueve desde el primer paso.}
\comprobacion el \cref{eje:5-grad} es un caso numérico de la segunda parte: $\partial\mathcal{L}/\partial B=(6,-6)\T\ne 0$ y $\partial\mathcal{L}/\partial A=0$.
\end{solucion}

% ---------------------------------------------------------------------
\section{Tema 6}

\begin{solucion}{eje:6-rag}
\planteamiento la regla del \cref{tema:6}: el ajuste cambia cómo responde el modelo; el contexto, qué sabe.
\pasos
\begin{enumerate}[label=\arabic*.]
  \item Coste de actualización: la normativa cambia cada mes. Con RAG basta con indexar los documentos nuevos; con ajuste fino habría que reentrenar cada mes, con riesgo de olvido (\cref{tema:3}).
  \item Fiabilidad: RAG pone el texto exacto de la norma en el contexto y permite citar la fuente; un modelo ajustado mezcla los datos en sus pesos y puede inventar detalles plausibles.
\end{enumerate}
\resultado \textbf{RAG}, por coste de actualización y por trazabilidad de las fuentes.
\comprobacion \cref{tab:eleccion}, fila «Conocimiento nuevo o cambiante».
\end{solucion}

\begin{solucion}{eje:6-cuantiles}
\datos $\mathcal{Q}=\{-1;\,-0{,}2770;\,0{,}2770;\,1\}$; $w=(0{,}9;\,-0{,}3;\,0{,}1;\,-1{,}8)$.
\planteamiento $c=\max\abs{w_j}$, normalizar, redondear al nivel más próximo y descuantizar (\cref{def:cuant-bloques}). Los puntos medios entre niveles son $\pm 0{,}6385$ y $0$.
\pasos
\begin{enumerate}[label=\arabic*.]
  \item $c=1{,}8$. Normalizados: $w/c=(0{,}5;\,-0{,}1667;\,0{,}05556;\,-1)$.
  \item Niveles más próximos: $0{,}5\to 0{,}2770$ (pues $0{,}5<0{,}6385$); $-0{,}1667\to-0{,}2770$; $0{,}05556\to 0{,}2770$; $-1\to-1$.
  \item Descuantizados: $1{,}8\cdot 0{,}2770=0{,}4986$, luego $\hat w=(0{,}4986;\,-0{,}4986;\,0{,}4986;\,-1{,}8)$.
  \item Errores $w-\hat w$: $(0{,}4014;\,0{,}1986;\,-0{,}3986;\,0)$.
  \item Norma: $\sqrt{0{,}1611+0{,}03944+0{,}1589+0}=\sqrt{0{,}3595}=0{,}5995$.
\end{enumerate}
\resultado \textbf{$\hat w=(0{,}4986;\,-0{,}4986;\,0{,}4986;\,-1{,}8)$; norma del error $0{,}5995$.} El peor parado en términos relativos es $w_3=0{,}1$, que se convierte en $0{,}4986$, casi cinco veces mayor: el libro de códigos no contiene el cero. NF4 incluye un nivel exactamente nulo (\cref{tema:6}), con el que $w_3$ tendría error $0{,}1$.
\comprobacion el máximo del bloque se representa sin error ($-1{,}8\to-1{,}8$), como ocurre siempre con la escala absmax.
\end{solucion}

\begin{solucion}{eje:6-doble}
\datos $N=65\cdot 10^{9}$; $b=4$; $B_1=64$; $B_2=256$.
\planteamiento \cref{prop:bits} y \cref{prop:doble}; bytes $=$ bits$/8$.
\pasos
\begin{enumerate}[label=\arabic*.]
  \item Sin doble cuantización: $4+32/64=4{,}5$ bits por parámetro; $65\cdot 10^{9}\cdot 4{,}5/8=36{,}56\cdot 10^{9}$ bytes $=36{,}56$~GB.
  \item Con doble cuantización: $4+8/64+32/(64\cdot 256)=4+0{,}125+0{,}001953=4{,}127$ bits; $65\cdot 10^{9}\cdot 4{,}127/8=33{,}53$~GB.
  \item Ahorro: $0{,}3730$ bits por parámetro; $65\cdot 10^{9}\cdot 0{,}3730/8=3{,}031$~GB.
\end{enumerate}
\resultado \textbf{$36{,}56$~GB sin doble cuantización, $33{,}53$~GB con ella: unos 3~GB de ahorro.}
\comprobacion $36{,}56-33{,}53=3{,}03$. En una GPU de 48~GB, esos 3~GB son una parte apreciable del margen para activaciones y adaptadores.
\end{solucion}

\begin{solucion}{eje:6-memoria}
\datos $N=7\cdot 10^{9}$; ajuste completo a 16~B por parámetro; QLoRA con base a $0{,}5$~B por parámetro y $N_{\mathrm{LoRA}}=2\cdot 10^{7}$ a 16~B.
\planteamiento \cref{prop:mem-qlora} con $b_{\mathrm{ef}}=4$.
\pasos
\begin{enumerate}[label=\arabic*.]
  \item (a) $7\cdot 10^{9}\cdot 16=112\cdot 10^{9}$~B $=112$~GB.
  \item (b) Base: $7\cdot 10^{9}\cdot 0{,}5=3{,}5$~GB. Adaptadores: $2\cdot 10^{7}\cdot 16=3{,}2\cdot 10^{8}$~B $=0{,}32$~GB. Total: $3{,}82$~GB.
  \item Razón: $112/3{,}82=29{,}32$.
\end{enumerate}
\resultado \textbf{(a) 112~GB; (b) $3{,}82$~GB; QLoRA usa unas $29{,}3$ veces menos memoria.}
\comprobacion los pesos solos pasan de 14~GB (bf16) a $3{,}5$~GB (un cuarto); el resto del ahorro procede de no guardar gradientes ni estado de Adam para $W$.
\end{solucion}

\begin{solucion}{eje:6-contabilidad}
\datos $V=50\,257$; $d=1600$; 48 capas; $d_{\mathrm{ff}}=4288$; $r=8$.
\planteamiento sumar por bloques; LoRA entrena $r(d+k)$ por matriz; memoria $=$ parámetros $\times$ bytes.
\pasos
\begin{enumerate}[label=\arabic*.]
  \item Atención por capa: $4d^{2}=4\cdot 2\,560\,000=10\,240\,000$.
  \item Red de alimentación por capa: $3\cdot 1600\cdot 4288=20\,582\,400$.
  \item RMSNorm por capa: $2\cdot 1600=3200$.
  \item Por capa: $30\,825\,600$; 48 capas: $1\,479\,628\,800$.
  \item \emph{Embeddings}: $2\cdot 50\,257\cdot 1600=160\,822\,400$. Norma final: $1600$.
  \item $N=1\,479\,628\,800+160\,822\,400+1600=1\,640\,452\,800\approx 1{,}64\cdot 10^{9}$.
  \item (b) LoRA por matriz: $8\cdot 3200=25\,600$. Total: $48\cdot 4\cdot 25\,600=4\,915\,200$, el $4\,915\,200/1\,640\,452\,800=0{,}2996\,\%$ de $N$.
  \item (c) Ajuste completo: $1\,640\,452\,800\cdot 16=26{,}25$~GB.
  \item LoRA: base $1\,640\,452\,800\cdot 2=3{,}281$~GB; adaptadores $4\,915\,200\cdot 16=0{,}07864$~GB; total $3{,}360$~GB.
  \item Razón: $26{,}25/3{,}360=7{,}813$.
\end{enumerate}
\resultado \textbf{(a) $N=1\,640\,452\,800$; (b) $4\,915\,200$ parámetros LoRA ($0{,}30\,\%$); (c) $26{,}25$~GB frente a $3{,}36$~GB, unas $7{,}8$ veces menos.}
\comprobacion LoRA equivale al $2r/d=16/1600=1\,\%$ de las matrices de atención: $48\cdot 10\,240\,000\cdot 0{,}01=4\,915\,200$.
\end{solucion}

\begin{solucion}{eje:6-rango-limite}
\datos matrices de $4096\times 4096$ y $1024\times 4096$.
\planteamiento comparar $r(d+k)$ con $dk$.
\pasos
\begin{enumerate}[label=\arabic*.]
  \item Como $d+k>0$, $r(d+k)<dk\iff r<dk/(d+k)$.
  \item $4096\times 4096$: $dk=16\,777\,216$, $d+k=8192$; límite $2048$. LoRA ahorra si $r<2048$.
  \item $1024\times 4096$: $dk=4\,194\,304$, $d+k=5120$; límite $819{,}2$. Como $r$ es entero, LoRA ahorra si $r\le 819$.
\end{enumerate}
\resultado \textbf{$r<dk/(d+k)$: $r<2048$ para $4096\times 4096$ y $r\le 819$ para $1024\times 4096$.}
\comprobacion con $d=k$ el límite es $d/2$, que da 2048 para $d=4096$. En la práctica se usan rangos de 4 a 64 porque la actualización útil tiene espectro concentrado (\cref{cor:error-lora}), el estado del optimizador queda mínimo y un rango pequeño regulariza.
\end{solucion}
